\documentclass{article}
\usepackage[margin=1in]{geometry}
\usepackage{mathtools, amsfonts, amsthm, xcolor, listings}

\title{HW06}
\author{Sam Ly}

\begin{document}
\maketitle

\section*{Exercise 1}
Collaborators: Sam Ly

\noindent
Tools and resources used: 
On Exercise 1, I just copied from the Assignment page. 
\pagebreak

\section*{Exercise 3}
\begin{itemize}
  \item Points attempted: 5
  \item Self-assessed score: 5/5.
  \item Time spent: 20 minutes.
  \item Difficulty: 3/10.
  \item Questions/Feedback: None.
\end{itemize}

Give examples of each of the following:
  
\begin{enumerate}
    \item {
        one group of order \(1\);

        The trivial group \(\{e\}\).
    }
    \item {
        one group of order \(p\) for each prime \(p\);

        The cyclic groups \(C_p\).
    }
    \item { 
        two groups of order \(4\) (both abelian);

        We have cyclic group \(C_4\) and the group of even integers \(2 \mathbb{Z} / 8\). 
    }

    \item {
        two groups of order \(6\) (one is abelian);

        We have the groups \(C_6\) and \(D_3\).
    }
    \item {
        five groups of order \(8\) (three are abelian, and one group we saw on Exam 1);

        We have the three abelian groups \(C_8\), \(C_4 \times C_2\), 
        \(C_2 \times C_2 \times C_2\), and the two non-abelian groups \(D_4\) 
        and the Quaternion group.
    }
    \item {
        two groups of order \(9\) (both are abelian); 

        We have the groups \(C_9\) and \(C_3 \times C_3\).
    }
    \item {
        two groups of order \(10\) (one is abelian).

        We have the the group \(C_10\) and \(D_5\).
    }
\end{enumerate}
\pagebreak
\section*{Choice Exercise 4}
\begin{itemize}
  \item Points attempted: 5
  \item Self-assessed score: 5/5.
  \item Time spent: 5 minutes.
  \item Difficulty: 1/10.
  \item Questions/Feedback: None.
\end{itemize}
\begin{enumerate}
    \item {
        (Do this part after you do the second part.) Write at least three sentences answering the following.

        \begin{enumerate}
            \item {
                What are a few ideas from HW 4 Exercise 2 that make a lot of sense to you?

                Part 1.b made a lot of sense to me, as I was able to relate it 
                back to a computer science concept I have learned previously.
            }

            \item {
                What ideas did you find most difficult?
        
                The most difficult problem for me was 2.d because I can `physically'
                see that the elements do not commute, but it is difficult to 
                form an intuition for why, and what this implies.
            }

            \item {
                Can you think of any ways to generalize the definitions or propositions that appeared in that problem?

                We can likely generalize the rule of all dihedral groups \(D_n\) 
                can be written in the form \(r^i f^j\) to some other types of 
                groups that have more than two `kinds' of elements. For example,
                \(a^i b^j c^k\), for some group generated by \(\{a,b,c\}\) with
                some relations \(a^n = e, b^3 = e, c^2 = e\). I am struggling 
                to really articulate this, but it `feels' like there is something
                there.
            }
        \end{enumerate}
    }

    \item {
    Copy/paste your answer from HW 4, re-read it, and make at least three changes that improve the exposition or correctness of the proofs.

    I changed two paragraphs of a.i, and one paragraph of a.iii.
    \begin{enumerate}
    \item {
        We defined the dihedral group in terms of its group presentation: 
        \[         D_n = \langle r, f \mid r^n = f^2 = (rf)^2 = e \rangle.       \]
      
        \begin{enumerate}
            \item {
                Prove that \(r^kf = fr^{n-k}\) for all \(0 < k < n\).

                \begin{proof}
                    First, we notice that \(rfrf = e\). By right muliplying by 
                    \(f\) we get \(rfr = f\). By left muliplying by \(r^{-1}\), 
                    we get \(fr = r^{-1}f\). Then, we left and right muliply by 
                    \(f\) on both sides to get \(rf = fr^{-1}\).
                    \emph{I fixed a type here.}

                    \emph{I added more reasoning to improve logical flow.}
                    We can continue this process for \(r^n f r^n f = e\) to show 
                    that \(r^1f = fr^{n-1}\). Now, to prove that 
                    \(r^kf = fr^{n-k}\) for all \(0 < k < n\), we must use 
                    induction. We have already proved the base case \(k = 1\). 
                    As our inductive hypothesis, we assume that the proposition 
                    is true for some value \(k\). Thus, \(r^kf = fr^{n-k}\).
                    Now, for by left multiplying by \(r\), we get 
                    \[r^{k+1}f = rfr^{n-k}.\]

                    Now, we apply \(rf = fr^{n-1}\) to the RHS, and get 
                    \[r^{k+1}f = fr^{-1}r^{n-k} = fr^{n- (k+1)}.\]
                    
                    Therefore, \(r^kf = fr^{n-k}\) for all \(0 < k < n\).
                \end{proof}
            }
            \item {
                Prove that every element can be written in the form \(r^kf^\ell\) 
                for \(0 \leq k < n\) and \(\ell \in \{0, 1\}\)

                \begin{proof}
                    We approach this problem from  a formal languages angle. 
                    We begin by constructing all possible `words' of the language 
                    formed by the alphabet \(\{r, s \}\). One example of such a 
                    word is \(rrrfrfffrffrrrrrrrf\). By definition, this language
                    must encompass all of \(D_n\) for any \(n\). 
                    
                    Then, we can form some 
                    basic substitution rules, in other words, a `grammar' that allows 
                    us to rewrite the word into a nicer form. By definition of the 
                    dihedral group, we have the rules: 
                    \begin{equation}
                        \underbrace{rr \dots r}_{m \text{ terms}} = r^m = r^{m-n}
                    \end{equation}
                    \begin{equation}
                        ff = e. \\
                    \end{equation}

                    Thus, we can rewrite our word into the form 
                    \(r^{k_0}f^{l_0} r^{k_1} f^{l_1} \dots \),
                    where \(0 \le k_i < n\) and \(l_i \in \{0, 1\}\). 

                    Now, from here, we apply the proposition proved above to 
                    combine and move all \(l_i\) to the right of the expression.

                    For example, if \(l_0 = 0\), then the expression becomes 
                    \(r^{k_0}r^{k_1} f^{l_1} \dots = r^{k_0 + k_1} f^{l_1} \dots \).
                    If \(l_0 = 1\), then the expression becomes 
                    \[r^{k_0}(fr^{k_1})f^{l_1} \dots = r^{k_0} (r^{n - k_1}f) f^{l_1} \dots \]
                    which we can simplify by combining \(ff^{l_1}\). 
                    This process can be repeated until we reach the form \(r^kf^\ell\) 
                    for \(0 \leq k < n\) and \(\ell \in \{0, 1\}\).
                \end{proof}
            }
            \item {
                Prove that the order (number of elements) of the group is \(D_n = 2n\).

                \begin{proof}
                    Now that we have established that all elements of \(D_n\) can 
                    be written as the form \(r^kf^\ell\) 
                    for \(0 \leq k < n\) and \(\ell \in \{0, 1\}\), this queestion
                    can be reframed as asking ``How many unique ways can we write 
                    \(r^kf^\ell\)?'' 
                    
                    \emph{I made this statement more detailed.}
                    Thus, for \(k\), we have \(n\) possible choices, and for \(\ell\),
                    we have 2 possible choices. Therefore, there are \(n * 2 = 2n\) 
                    total choices possible. Thus, \(|D_n| = 2n\).
                \end{proof}
            }

        \end{enumerate}
    }

    \item {
        We say that a group \((G,*)\) is an abelian group if and only if 
        \(a * b = b * a\) for all \(a, b \in G\).
      
        \begin{enumerate}
            \item {
                Prove that if \(G\) is generated by \(g_1\) and \(g_2\) then 
                \(G\) is abelian if and only if \(g_1 * g_2 = g_2 * g_1\).


                \begin{proof}
                    First, we prove that if \(a * b = b * a\) for all \(a, b \in G\),
                    then \(g_1 * g_2 = g_2 * g_1\). By setting \(a = g_1\) and 
                    \(b = g_2\),  then \(g_1 * g_2 = g_2 * g_1\).

                    The converse ``if \(g_1 * g_2 = g_2 * g_1\), then \(a * b = b * a\) for all \(a, b \in G\)''
                    is much more difficult to prove. 

                    Notice that 
                    \begin{align}
                        a &= a_1 * a_2 * \dots * a_p = \prod_{i=1}^p a_i \\
                        b &= b_1 * b_2 * \dots * b_q = \prod_{i=1}^q b_i \\
                        c &= a * b = a_1 * a_2 * \dots * a_p * b_1 * b_2 * \dots * b_q
                    \end{align}
                    where \(a_i, b_i \in \{g_1, g_2\}\).

                    Notice that because \(g_1 * g_2 = g_2 * g_1\), we can swap 
                    any two adjacent terms of \(c\).

                    Therefore, we can rewrite \(c\) as 
                    \begin{align}
                        c &= a * b = b_1 * b_2 * \dots * b_q * a_1 * a_2 * \dots * a_p.
                    \end{align}

                    Therefore, \(G\) is generated by \(g_1\) and \(g_2\) then 
                    \(G\) is abelian if and only if \(g_1 * g_2 = g_2 * g_1\).
                \end{proof}
            }
        
            \item {
                Suppose \(G\) is generated by \(g_1, g_2, \dots, g_n\). Do you 
                think that \(G\) is abelian if and only if \(g_i * g_j = g_j * g_i\) 
                for all \(i, j \in \{1, 2, \dots, n\}\)? Explain your reasoning, 
                but you do not have to provide a proof.

                I do think so. We can likely apply some inductive proof to show that 
                adding a new \(g_k\) that is mutually commutative with all other 
                \(g_i\) does not impact the `abelian-ness' of the overall group 
                \(G\).
            }
        
            \item {
                Prove that cyclic groups \(C_n\) are abelian.
        
                \begin{proof}
                    By definition, 
                    \[C_n = \langle r \mid r^n = e \rangle.\]

                    Because the cyclic groups \(C_n\) are defined by a generator 
                    with one non-identity element, then \(C_n\) must be abelian
                    by associativity.

                    \begin{align}
                        r^k * r^l &= (\underbrace{r * r * \dots * r}_{k}) * (\underbrace{r * r * \dots * r}_{l} )\\
                                  &= (\underbrace{r * r * \dots * r}_{l}) * (\underbrace{r * r * \dots * r}_{k}) \\
                                  &= r^l * r^k.
                    \end{align}
                \end{proof}
            }

            \item {
                Prove that the general linear group of \(2 \times 2\) 
                real-valued matrices 
                \[GL_2(\mathbb{R}) = \{A \in \mathbb{R}^{n \times n} \mid \det(A) \neq 0\}\] 
                is not abelian.
        
                \begin{proof}
                    Let 
                    \[
                        A = 
                            \begin{bmatrix}
                                a & b \\
                                c & d 
                            \end{bmatrix},
                        B = 
                            \begin{bmatrix}
                                p & q \\
                                r & s
                            \end{bmatrix}.
                    \]

                    Then, \(A \times B = \begin{bmatrix}
                        ap + br & aq + bs \\
                        cp + dr & cq + ds 
                    \end{bmatrix} \) and 
                    \(B \times A = \begin{bmatrix}
                        ap + cq & bp + dq \\
                        ar + cs & br + ds 
                    \end{bmatrix}\).

                    In general, these are not the same matrix.
                \end{proof}
            }
            
            \item {
                Prove that the dihedral group \(D_n\) is abelian if and only if \(n \leq 2\).
                
                \begin{proof}
                    First, by definition, the smallest dihedral group is \(D_1\), 
                    which is just the identity. This trivial group must be abelian. 

                    Then, for \(n = 2\),
                    \[D_2 = \langle r, f \mid r^2 = f^2 = e \rangle. \]
                    Notice that in this case, \(r = f\), therefore, \(D_2\) is 
                    abelian.

                    Then, to prove the ``other direction'', we must prove that 
                    a dihedral group \(D_n\) being abelian implies that \(n \leq 2\).
                    We can prove this via contrapositive. Thus, we for a dihedral
                    group \(D_n\) where \(n > 2\), \(D_n\) is not abelian.

                    Suppose we have \(D_n\) where \(n > 2\). We choose the element
                    \(rf \in D_n\). Now, \(rf = fr^{n-1} \not= fr\), so \(D_n\) is 
                    not abelian. 

                    Therefore, \(D_n\) is abelian if and only if \(n \leq 2\).
                \end{proof}
            }
        \end{enumerate}
    }

    \item {
        The direct product of groups is a way of combining two groups into a new group. The notation can look a bit intimidating, but the basic idea is simple: just treat the groups independently. Here's the formal definition, but keep in mind that \(*\), \(\star\), and \(\odot\) are just names for binary operations like \(+\), \(\times\), and \(\circ\).
      
        Given two groups \((G_1,*)\) and \((G_2,\star)\), we define the direct product of groups as \[           G_1 \times G_2 = \{(g_1, g_2) \mid g_1 \in G_1 \text{ and } g_2 \in G_2\}         \] with the binary operation, \(\odot\), where \[           (g_1, g_2)\odot(g'_1,g'_2) = (g_1 * g'_1, g_2 \star g'_2).         \]
      
        \begin{enumerate}
            \item {
                Let \({(G_1, *) = (D_4, *)}\) be the dihedral group of the square, 
                and let \({(G_2, \star) = (S_3, \circ)}\) be the symmetric group 
                under function composition.  Show that in \(G_1 \times G_2\), 
                \[\left(f, (132) \right) \odot \left(r, (23) \right) = \left(r^3f, (13) \right).\]

                \begin{align}
                    \left(f, (132) \right) \odot \left(r, (23) \right) 
                        &= (fr, (132) \circ (23)) \\
                        & = (r^3f, (2)(31)) = (r^3f, (13)).
                \end{align}
            }
            
            \item {
                Prove that for any two groups \(G_1\) and \(G_2\), the direct product of groups is a group.
                
                \begin{proof}
                    Let \(G = G_1 \times G_2\). We verify that \(G\) follows the group 
                    axioms:
                    \begin{enumerate}
                        \item[Closure:] {
                            Let \(g_1, g_1' \in G_1\) and \(g_2, g_2' \in G_2\).
                            since \(g_1g_1' \in G_1\), and \(g_2g_2' \in G_2\),
                            \((g_1g_1', g_2,g_2') \in G\).
                        }
                        \item[Identity:] {
                            Let \(e_1, e_2\) be the identity elements of \(G_1\) 
                            and \(G_2\) respectively. \((e_1, e_2)\) is the identity
                            of \(G\).
                        }
                        \item[Associativity:] {
                            Since the operations on the first and second elements of 
                            \(G\) are both associative, associativity is inherited 
                            to the direct product.
                        }
                        \item[Inverse:] {
                            Similarly, because the left and right elements each 
                            independently have an inverse, the tuple containing 
                            both will be the equivalent inverse in \(G\).
                        }
                    \end{enumerate}
                \end{proof}
            }

            \item {
                Prove that if \(G_1 \times G_2\) is abelian if and only if \(G_1\) and \(G_2\) are both abelian.
        
                First, if \(G = G_1 \times G_2\) is abelian, then for 
                \((a,b), (a', b') \in G\),
                
                \begin{align}
                    (a,b) \odot (a',b') &= (a', b') \odot (a, b) \\
                    (a*a', b \star b')  &= (a' * a, b' \star b) 
                \end{align}
                
                Since \(a*a' = a'*a\) and \(b \star b' = b'\star b\), \(G_1\) and 
                \(G_2\) are abelian.

                Now, given \(G_1\) and \(G_2\) are abelian, reverse the above 
                algebra to see that \(G\) must also be abelian. Thus, 
                \(G_1 \times G_2\) is abelian if and only if \(G_1\) and \(G_2\) 
                are both abelian.

            }

            \item {
                Let \(\mathbb{Z}/k\mathbb{Z}\) denote the group of the integers mod \(k\) under addition. Let \(G = \mathbb{Z}/2\mathbb{Z} \times \mathbb{Z}/3\mathbb{Z}\). Compute the order of \(G\). Prove that \(G\) has an element of order \(|G|\).
        
                The order of a direct product of two groups is the product of the 
                orders. Thus,
                \[|G| = |\mathbb{Z}/2\mathbb{Z}| \times |\mathbb{Z}/3\mathbb{Z}| = 2 \times 3 = 6.\]

                The element with order 6 is \((1,1)\).
            }

            \item {
                Let \(G = \mathbb{Z}/2\mathbb{Z} \times \mathbb{Z}/2\mathbb{Z}\). 
                Compute the order of \(G\). Prove the \(G\) has no element of order \(|G|\).

                Similarly, \(|G| = 2 \times 2 = 4\). Now, the elements of \(G\) 
                are \(\{ (0,0), (0,1), (1,0), (1,1)  \}\). The orders of all of these 
                elements is \(2\).
            }
        \end{enumerate}
    }
\end{enumerate}
\
    }
\end{enumerate}
\pagebreak


\section*{Choice Exercise 7 [10]}
\begin{itemize}
  \item Points attempted: 10
  \item Self-assessed score: 10/10.
  \item Time spent: 30 minutes.
  \item Difficulty: 3/10.
  \item Questions/Feedback: None.
\end{itemize}

\begin{enumerate}
    \item {
        Let \(G\) be a group with order \(n\), and denote the order of an element \(g \in G\) by \(|g|\). Then the order of \(G\) is equal to \[         \gcd(|g_1|, |g_2|, \dots, |g_n|) = n.       \]

        False. The order of group \(G\) should be at least the lcm of the orders 
        of the elements.

        Answer: False. My reasoning was maybe correct.
    }

    \item {
        If \(\phi\colon G \to H\) is a homomorphism that is 1-to-1, then \(\phi\colon G \to \operatorname{Im}(\phi)\) is a isomorphism.
    
        True.

        Answer: Similar reasoning.
    }

    \item {
        If \(\phi\colon G \to H\) is a homomorphism that is onto, then \(\operatorname{Im}(\phi) \triangleleft H\).
    
        False. Consider the case where \(G\) is a subgroup of \(H\), but not 
        necessarily a normal subgroup. Now, let \(\phi(g) = g\) for \(g \in G\). 

        Answer: True. I overlooked the fact that the homomorphism was onto.
    }

    \item {
        If \(G\) is a non-abelian group and \(H\) is abelian, then the only homomorphism \(\phi \colon G \to H\) is the trivial homomorphism defined by \(\phi(g) = e_H\) for all \(g \in G\).

        False. Consider the homomorphism from \(\phi \colon D_3 \to C_2\), where 
        \(\phi(r^n) = e\) and \(\phi(f) = 1\). 

        Answer: False. I came to the exact same conclusion, with the same counterexample.
    }

    \item {
        If \(\varphi \colon G \to H\) is an isomorphism, then there always exists an isomorphism \({\psi \colon H \to G}\).
    
        True.

        Answer: True, similar reasoning.
    }

    \item {
        Let \(G \times H\) be a direct product of groups. Then the set 
        \[         N = \{(g, e_H) \mid g \in G\} \subseteq G \times H       \] 
        is a normal subgroup of \(G \times H\).
    
        False. This is only a normal subgroup if \(H\) is abelian.

        Answer: True. I reasoned incorrectly. If I were to correct my reasoning, 
        I would note that \(h * e_H = e_H * h = h\).
    }

    \item {
        If \(G/N \cong H\) then \(G \cong N \times H\).
    
        True. Under the isomorphism \(\phi\), each of the cosets \(gN\) can be 
        mapped to a unique \(h \in H\). By Langrange's theory, there must be 
        \(|N|\) unique \(g \in G\) that maps to each \(h\). Also, \(gN\) is 
        closed under \(n \in N\). Therefore, we can construct an isomorphism from 
        \(G\) to \(N \times H\).

        Answer: False, my reasoning onl showed that the orders of the groups were 
        compatible, not that an isomorphism actually exists.
    }

    \item {
        The map \(\varphi\colon \mathbb{Z}/8\mathbb{Z} \to C_{16}\) given by \(\varphi(1) = r \in C_{16}\) is a homomorphism.
    
        False.

        Answer: Correct.
    }

    \item {
        There is a homomorphism \(\psi\colon D_6 \to C_6\) such that \(\psi(r) = r \in C_6\) and \(\psi(f) = e \in C_6\).
    
        False. In \(D_6\), \(f\) can be thought of as ``inverting'' all subsequent 
        \(r\) elements, so this is not a homomorphism.

        Answer: Correct, similar reasoning.
    }

    \item {
        There is a homomorphism \(\phi\colon C_6 \to D_6\) such that \(\phi(r) = r \in D_6\).
    
        True.

        Answer: Correct, but I didn't check the fact that \(\phi\) was structure 
        preserving with the relations. I just wrote out the possibilities.
    }
\end{enumerate}


\end{document}
